% !TEX root = ../main.tex
%% Chapitre 5 — Nombres réels et suites numériques
%% Programme officiel MPSI/MP2I (BO spécial n°1 du 11/02/2021), premier semestre,
%% section « Nombres réels et suites numériques » : a) ensembles de nombres usuels ;
%% b) borne supérieure ; c) généralités sur les suites ; d) limites ; e) suites
%% monotones ; f) suites extraites et Bolzano-Weierstrass ; g) traduction
%% séquentielle ; h) suites complexes ; i) suites particulières.
%% Résultats vérifiés par verifications/chapitre_05.py.

\begin{ouverture}[suites]{5}{Nombres réels et suites numériques}{Partie 1 · Premier semestre}
  {Borne supérieure\enspace·\enspace Limites\enspace·\enspace Suites monotones\enspace·\enspace
   Suites adjacentes\enspace·\enspace Suites extraites\enspace·\enspace Suites récurrentes}
  \objectif{Manipuler les ensembles de nombres, la densité et les bornes supérieures et inférieures.}
  \objectif{Maîtriser la définition de la limite d'une suite et les théorèmes d'opérations.}
  \objectif{Utiliser encadrements, comparaisons, limite monotone et suites adjacentes.}
  \objectif{Raisonner avec les suites extraites, jusqu'au théorème de Bolzano-Weierstrass.}
  \objectif{Résoudre les récurrences linéaires d'ordre 2 et les récurrences arithmético-géométriques.}
  \objectif{Étudier une suite définie par $u_{n+1}=f(u_n)$ : intervalle stable, monotonie, limite.}
  \prerequis{Logique et récurrences (chapitre 1) ; inégalités, valeur absolue, partie entière (chapitre 2).}
  \prerequis{Suites arithmétiques et géométriques du lycée.}
  \prerequis{Fonctions usuelles et inégalités de convexité élémentaires (chapitre 4).}
\end{ouverture}

%% =====================================================================
%%  L'ESSENTIEL
%% =====================================================================
\essentiel

\essSection{1}{Nombres réels, décimaux, densité}

\begin{definition}[Les ensembles de nombres]
$\N\subset\Z\subset\mathbb{D}\subset\Q\subset\R$, où $\mathbb{D}$ désigne
l'ensemble des \textbf{décimaux} (les $\frac{a}{10^{n}}$, $a\in\Z$,
$n\in\N$) et $\Q$ celui des \textbf{rationnels}. Un réel non rationnel est
dit \textbf{irrationnel} ($\sqrt2$, $\mathrm{e}$, $\pi$\dots). Les
constructions de ces ensembles sont hors programme. Pour tout réel $x$ et
tout $n\in\N$, l'encadrement
\[
  \frac{\lfloor10^{n}x\rfloor}{10^{n}}\;\leq\;x\;<\;\frac{\lfloor10^{n}x\rfloor+1}{10^{n}}
\]
fournit les \textbf{approximations décimales} de $x$ à $10^{-n}$ près, par
défaut et par excès.
\end{definition}

\begin{theoreme}[Densité]
Tout intervalle ouvert non vide de $\R$ contient au moins un rationnel et un
irrationnel : on dit que $\Q$ et $\R\setminus\Q$ sont \textbf{denses} dans
$\R$. Plus généralement, une partie de $\R$ est dense dans $\R$ si elle
rencontre tout intervalle ouvert non vide ; c'est le cas de $\mathbb{D}$.
\end{theoreme}

\begin{definition}[Droite achevée]
$\overline{\R}=\R\cup\ens{-\infty,+\infty}$, avec les conventions d'ordre et
de calcul usuelles ; elle permet d'énoncer uniformément les résultats sur les
limites.
\end{definition}

\essSection{2}{La propriété de la borne supérieure}

\noindent
\begin{minipage}[c]{0.56\linewidth}
\begin{definition}[Borne supérieure]
Soit $X$ une partie non vide de $\R$. La \textbf{borne supérieure}
$\sup X$ est, lorsqu'elle existe, \emph{le plus petit des majorants} de $X$ ;
la borne inférieure $\inf X$ est le plus grand des minorants.
Si $X$ possède un plus grand élément, alors $\sup X=\max X$ ; mais la borne
supérieure peut très bien ne pas être atteinte :
$\sup\intFO{0}{1}=1\notin\intFO{0}{1}$.
\end{definition}
\end{minipage}\hfill
\begin{minipage}[c]{0.4\linewidth}
\centering
\begin{tikzpicture}[line cap=round,xscale=1.35]
  \draw[otMuted,line width=0.5pt,-{Stealth[length=1.8mm]}] (-0.3,0) -- (4.3,0);
  \draw[otNavy,line width=2.4pt] (0.4,0) -- (1.3,0);
  \draw[otNavy,line width=2.4pt] (1.7,0) -- (2.6,0);
  \node[below,font=\scriptsize,text=otNavy] at (1.5,-0.12) {$X$};
  \draw[otGreenDark,line width=1.1pt] (2.6,-0.09) -- (2.6,0.09);
  \node[above,font=\scriptsize,text=otGreenDark] at (2.6,0.12) {$\sup X$};
  \foreach \x in {3.0,3.4,3.8}{\draw[otMuted,line width=0.7pt] (\x,-0.07) -- (\x,0.07);}
  \node[above,font=\scriptsize,text=otMuted] at (3.45,0.12) {majorants};
\end{tikzpicture}\\[0.5mm]
\legende{$\sup X$ : le plus petit des majorants de $X$.}
\end{minipage}

\begin{theoreme}[Propriété fondamentale de $\R$]
Toute partie de $\R$ \emph{non vide et majorée} admet une borne supérieure ;
toute partie non vide et minorée admet une borne inférieure. C'est cette
propriété qui distingue $\R$ de $\Q$ : l'ensemble
$\ens{x\in\Q \mid x^{2}<2}$ n'a pas de borne supérieure \emph{dans} $\Q$.
\end{theoreme}

\begin{savoirfaire}[Caractériser une borne supérieure]
$M=\sup X$ si et seulement si :
\begin{itemize}[label={\color{otGreenDark}\textbullet},leftmargin=1.4em,itemsep=0.6mm,topsep=0.5mm]
  \item $M$ majore $X$ : pour tout $x\in X$, $x\leq M$ ;
  \item rien de plus petit ne majore : pour tout $\varepsilon>0$, il existe
    $x\in X$ tel que $x>M-\varepsilon$.
\end{itemize}
\emph{Traduction séquentielle} : si $X$ est non vide et majorée, il existe une
suite d'éléments de $X$ qui converge vers $\sup X$ (et une suite tendant vers
$+\infty$ si $X$ n'est pas majorée).
\end{savoirfaire}

\begin{propriete}[Caractérisation des intervalles]
Une partie $X$ de $\R$ est un intervalle si et seulement si, pour tous
$a,b\in X$ avec $a\leq b$, on a $\intFF{a}{b}\subset X$ (partie
\emph{convexe} de $\R$).
\end{propriete}

\essSection{3}{Suites réelles : vocabulaire}

\begin{definition}
Une suite réelle $(u_n)_{n\in\N}$ est \textbf{majorée} (resp. minorée,
bornée) si l'ensemble de ses termes l'est ; elle est bornée si et seulement
si $(|u_n|)$ est majorée. Elle est \textbf{croissante} si
$u_{n+1}\geq u_n$ pour tout $n$, \textbf{stationnaire} si elle est constante
à partir d'un certain rang. Une suite peut être définie de façon
\emph{explicite} ($u_n=f(n)$), \emph{implicite} ($u_n$ unique solution d'une
équation dépendant de $n$) ou \emph{par récurrence}.
\end{definition}

\essSection{4}{Limite d'une suite}

\begin{definition}[Limite, avec des inégalités larges]
$(u_n)$ \textbf{converge} vers $\ell\in\R$ si :
\[
  \forall\varepsilon>0,\ \exists N\in\N,\ \forall n\geq N,\quad|u_n-\ell|\leq\varepsilon.
\]
On note $u_n\longrightarrow\ell$ ou $\lim u_n=\ell$ ; la limite, si elle
existe, est \emph{unique}. De même, $u_n\to+\infty$ si tout
$\intFO{A}{+\infty}$ contient les $u_n$ à partir d'un certain rang. Une
suite qui ne converge pas (vers un réel) est dite \textbf{divergente}.
\end{definition}

\begin{theoreme}[Premiers résultats]
Toute suite convergente est bornée. Les limites se combinent : combinaison
linéaire, produit, quotient (dénominateur de limite non nulle). Le produit
d'une suite \emph{bornée} par une suite de limite \emph{nulle} tend vers $0$.
Par passage à la limite, les inégalités \emph{larges} se conservent ; et si
$u_n\to\ell>0$, alors $u_n>0$ à partir d'un certain rang.
\end{theoreme}

\noindent
\begin{minipage}[c]{0.52\linewidth}
\begin{theoreme}[Existence de limites par comparaison]
\begin{itemize}[label={\color{otNavy}\textbullet},leftmargin=1.4em,itemsep=0.6mm,topsep=0.8mm]
  \item \etape{Encadrement} si $a_n\leq u_n\leq b_n$ à partir d'un certain
    rang et si $a_n\to\ell$ et $b_n\to\ell$, alors $u_n\to\ell$.
  \item \etape{Minoration} si $a_n\leq u_n$ et $a_n\to+\infty$, alors
    $u_n\to+\infty$ (et l'analogue en $-\infty$).
  \item \etape{Majoration de l'écart} si $|u_n-\ell|\leq v_n$ avec
    $v_n\to0$, alors $u_n\to\ell$.
\end{itemize}
\end{theoreme}
\end{minipage}\hfill
\begin{minipage}[c]{0.44\linewidth}
\centering
\begin{tikzpicture}[line cap=round]
  \draw[otMuted,line width=0.5pt,-{Stealth[length=1.6mm]}] (-0.2,0) -- (4.6,0) node[below,font=\scriptsize] {$n$};
  \draw[otNavy,dashed,line width=0.5pt] (-0.1,1.1) -- (4.5,1.1);
  \node[left,font=\scriptsize,text=otNavy] at (-0.1,1.1) {$\ell$};
  \foreach \n in {1,...,13}{
    \fill[otRule!70!otMuted] ({0.32*\n},{1.1+0.9/\n}) circle (0.85pt);
    \fill[otRule!70!otMuted] ({0.32*\n},{1.1-0.9/\n}) circle (0.85pt);
    \fill[otGreenDark] ({0.32*\n},{1.1+0.82*sin(137*\n)/\n}) circle (1.1pt);}
  \node[right,font=\scriptsize,text=otMuted] at (4.3,1.45) {$b_n$};
  \node[right,font=\scriptsize,text=otMuted] at (4.3,0.78) {$a_n$};
  \node[right,font=\scriptsize,text=otGreenDark] at (4.3,1.12) {$u_n$};
\end{tikzpicture}\\[0.5mm]
\legende{Encadrement : $(u_n)$, prise en étau, converge vers $\ell$.}
\end{minipage}

\essSection{5}{Suites monotones, suites adjacentes}

\begin{theoreme}[Limite monotone]
Toute suite croissante et majorée converge, vers la borne supérieure de ses
termes ; toute suite croissante non majorée tend vers $+\infty$. Une suite
croissante possède donc toujours une limite dans $\overline{\R}$ (énoncés
analogues en décroissant).
\end{theoreme}

\noindent
\begin{minipage}[c]{0.52\linewidth}
\begin{theoreme}[Suites adjacentes]
Deux suites $(a_n)$ et $(b_n)$ sont \textbf{adjacentes} si $(a_n)$ est
croissante, $(b_n)$ décroissante, et $b_n-a_n\to0$. Deux suites adjacentes
convergent, vers une \emph{même} limite $\ell$, et pour tout $n$ :
\[
  a_n\leq\ell\leq b_n :
\]
chaque couple $(a_n,b_n)$ fournit un encadrement de $\ell$.
\end{theoreme}
\end{minipage}\hfill
\begin{minipage}[c]{0.44\linewidth}
\centering
\begin{tikzpicture}[line cap=round]
  \draw[otMuted,line width=0.5pt,-{Stealth[length=1.6mm]}] (-0.2,0) -- (4.6,0) node[below,font=\scriptsize] {$n$};
  \draw[otNavy,dashed,line width=0.5pt] (-0.1,1.1) -- (4.5,1.1);
  \node[left,font=\scriptsize,text=otNavy] at (-0.1,1.1) {$\ell$};
  \foreach \n in {1,...,12}{
    \fill[otGreenDark] ({0.35*\n},{1.1-1.05/\n}) circle (1.1pt);
    \fill[otOrange] ({0.35*\n},{1.1+0.8/\n}) circle (1.1pt);}
  \node[right,font=\scriptsize,text=otOrange!80!black] at (4.3,1.32) {$b_n$};
  \node[right,font=\scriptsize,text=otGreenDark] at (4.3,0.85) {$a_n$};
\end{tikzpicture}\\[0.5mm]
\legende{Adjacentes : $(a_n)$ monte, $(b_n)$ descend, l'écart fond — la limite est coincée.}
\end{minipage}

\essSection{6}{Suites extraites, théorème de Bolzano-Weierstrass}

\begin{definition}[Suite extraite]
Une \textbf{suite extraite} (ou sous-suite) de $(u_n)$ est une suite
$(u_{\varphi(n)})$, où $\varphi:\N\to\N$ est strictement croissante.
Si $(u_n)$ possède une limite (finie ou infinie), toutes ses suites
extraites possèdent la même limite.
\end{definition}

\begin{propriete}[Deux usages essentiels]
\begin{itemize}[label={\color{otNavy}\textbullet},leftmargin=1.4em,itemsep=0.6mm,topsep=0.8mm]
  \item \etape{Divergence} si deux suites extraites ont des limites
    différentes, la suite diverge : ainsi $\bigl((-1)^{n}\bigr)$ n'a pas de
    limite.
  \item \etape{Recollement} si $(u_{2n})$ et $(u_{2n+1})$ convergent vers la
    \emph{même} limite $\ell$, alors $(u_n)$ converge vers $\ell$.
\end{itemize}
\end{propriete}

\begin{theoreme}[Bolzano-Weierstrass]
De toute suite réelle \emph{bornée}, on peut extraire une suite convergente.
Le principe de la démonstration est la \emph{dichotomie} : on coupe en deux
un segment contenant une infinité de termes, on garde une moitié qui en
contient encore une infinité, et on recommence.
\end{theoreme}

\begin{definition}[Suites complexes]
Les définitions (suite bornée, convergence avec $|z_n-\ell|$) s'étendent aux
suites complexes : $(z_n)$ converge vers $\ell$ si et seulement si
$\bigl(\operatorname{Re}z_n\bigr)$ et $\bigl(\operatorname{Im}z_n\bigr)$
convergent vers $\operatorname{Re}\ell$ et $\operatorname{Im}\ell$.
Le théorème de Bolzano-Weierstrass subsiste pour les suites complexes
bornées.
\end{definition}

\essSection{7}{Suites particulières}

\begin{propriete}[Arithmético-géométriques]
Pour $u_{n+1}=a\,u_n+b$ avec $a\in\Cx\setminus\ens{1}$ : on cherche la
\emph{constante} $\omega=a\omega+b$, soit $\omega=\dfrac{b}{1-a}$ ; alors
$(u_n-\omega)$ est géométrique de raison $a$, d'où
$u_n=\omega+(u_0-\omega)\,a^{n}$. La suite converge (vers $\omega$) si
$|a|<1$ ou si $u_0=\omega$.
\end{propriete}

\begin{theoreme}[Récurrences linéaires d'ordre 2]
Pour $u_{n+2}=a\,u_{n+1}+b\,u_n$ ($a,b$ réels, $b\neq0$), on introduit
l'\textbf{équation caractéristique} $r^{2}=ar+b$.
\begin{itemize}[label={\color{otNavy}\textbullet},leftmargin=1.4em,itemsep=0.6mm,topsep=0.8mm]
  \item Deux racines réelles $r_1\neq r_2$ : $u_n=A\,r_1^{n}+B\,r_2^{n}$ ;
  \item racine double $r_0\neq0$ : $u_n=(A+Bn)\,r_0^{n}$ ;
  \item racines complexes conjuguées $r\,\mathrm{e}^{\pm\mathrm{i}\theta}$ :
    $u_n=r^{n}\bigl(A\cos(n\theta)+B\sin(n\theta)\bigr)$ (suites réelles).
\end{itemize}
Les constantes $A$, $B$ se déterminent avec $u_0$ et $u_1$.
\end{theoreme}

\noindent
\begin{minipage}[c]{0.52\linewidth}
\begin{savoirfaire}[Étudier $u_{n+1}=f(u_n)$]
\begin{enumerate}[label=\otLabel{\arabic*.},leftmargin=1.6em,itemsep=0.6mm,topsep=0.5mm]
  \item Trouver un \emph{intervalle stable} $I$ ($f(I)\subset I$) contenant
    $u_0$ : alors tous les $u_n$ restent dans $I$ (récurrence).
  \item Étudier la monotonie, souvent via le \emph{signe de} $f(x)-x$ sur
    $I$ (ou la croissance de $f$).
  \item Conclure par limite monotone, puis identifier la limite : si
    $u_n\to\ell$ et $f$ continue en $\ell$, alors $f(\ell)=\ell$.
\end{enumerate}
\end{savoirfaire}
\end{minipage}\hfill
\begin{minipage}[c]{0.44\linewidth}
\centering
\begin{tikzpicture}[scale=1.5,line cap=round]
  \draw[otMuted,line width=0.5pt,-{Stealth[length=1.6mm]}] (-0.15,0) -- (2.35,0) node[below,font=\scriptsize] {$x$};
  \draw[otMuted,line width=0.5pt,-{Stealth[length=1.6mm]}] (0,-0.15) -- (0,2.3);
  \draw[otRule!80!otMuted,line width=0.6pt] (0,0) -- (2.2,2.2);
  \draw[otNavy,line width=1pt] plot[domain=0:2.2,samples=50] ({\x},{sqrt(2+\x)});
  \draw[otGreen,line width=0.8pt]
    (0,0) -- (0,1.414) -- (1.414,1.414) -- (1.414,1.848) -- (1.848,1.848)
    -- (1.848,1.962) -- (1.962,1.962);
  \fill[otNavy] (2,2) circle (1.1pt);
  \foreach \x in {0,1.414,1.848}{\fill[otGreenDark] (\x,0) circle (0.9pt);}
  \draw[otGreen!70!otMuted,densely dotted,line width=0.5pt] (1.414,1.414) -- (1.414,0);
  \draw[otGreen!70!otMuted,densely dotted,line width=0.5pt] (1.848,1.848) -- (1.848,0);
  \node[below,font=\tiny,text=otGreenDark] at (0,-0.04) {$u_0$};
  \node[below,font=\tiny,text=otGreenDark] at (1.414,-0.04) {$u_1$};
  \node[below,font=\tiny,text=otGreenDark] at (1.9,-0.04) {$u_2$};
  \node[font=\scriptsize,text=otNavy] at (0.55,1.62) {$y=f(x)$};
  \node[font=\scriptsize,text=otMuted,rotate=45] at (1.05,0.88) {$y=x$};
\end{tikzpicture}\\[0.5mm]
\legende{La « toile d'araignée » : les termes montent vers le point fixe.}
\end{minipage}

\begin{attention}[Pièges fréquents]
\begin{itemize}[label={\color{otOrange}\textbullet},leftmargin=1.4em,itemsep=0.6mm,topsep=0.5mm]
  \item $\sup X$ n'est pas toujours un élément de $X$ : ne pas confondre borne
    supérieure et plus grand élément.
  \item Une suite croissante ne converge que si elle est \emph{majorée} ; une
    suite bornée peut diverger ($(-1)^{n}$).
  \item Le passage à la limite transforme les inégalités \emph{strictes} en
    inégalités \emph{larges} : $\frac1n>0$ mais $\lim\frac1n=0$.
  \item $f(\ell)=\ell$ ne prouve pas la convergence : il faut \emph{d'abord}
    établir que la suite converge.
  \item Pour les suites adjacentes, vérifier les \emph{trois} hypothèses —
    l'écart qui tend vers $0$ ne suffit pas.
  \item Une suite extraite qui converge ne dit rien de la suite entière :
    $(-1)^{n}$ a des sous-suites convergentes.
\end{itemize}
\end{attention}

%% =====================================================================
%%  EXERCICES
%% =====================================================================
\exercices
\rubrique{Applications}

\begin{exercice}[1]{Rationnels et irrationnels}
\begin{questions}
  \item Vrai ou faux (justifier ou donner un contre-exemple) : la somme de
    deux irrationnels est irrationnelle ; le produit de deux irrationnels
    est irrationnel.
  \item Montrer que $\sqrt2+\sqrt3$ est irrationnel.
    \emph{On pourra élever au carré et utiliser l'irrationalité de
    $\sqrt6$, que l'on admettra.}
  \item Donner les approximations décimales de $\sqrt2$ à $10^{-3}$ près,
    par défaut et par excès.
\end{questions}
\end{exercice}

\begin{exercice}[1]{Bornes supérieures et inférieures}
Pour chacun des ensembles suivants, déterminer, s'ils existent, la borne
supérieure, la borne inférieure, le plus grand et le plus petit élément :
\[
  A=\intFO{0}{1},\qquad
  B=\Bigl\{\frac1n \mid n\in\N^{*}\Bigr\},\qquad
  C=\Bigl\{(-1)^{n}+\frac1n \mid n\in\N^{*}\Bigr\},\qquad
  D=\ens{x\in\Q \mid x^{2}<2}.
\]
\end{exercice}

\begin{exercice}[1]{Majorée, minorée, bornée ?}
Étudier si les suites suivantes sont majorées, minorées, bornées :
\begin{questions}
  \item $u_n=2-\dfrac{1}{n+1}$ ;
  \item $v_n=(-1)^{n}\,n$ ;
  \item $w_n=\sin n+\dfrac1n$\quad($n\geq1$) ;
  \item $t_n=n^{2}-n$.
\end{questions}
\end{exercice}

\begin{exercice}[1]{Monotonie}
Étudier la monotonie des suites suivantes (définies pour $n\geq1$) :
\begin{questions}
  \item $u_n=n^{2}-n$ ;
  \item $v_n=\dfrac{2^{n}}{n!}$\quad\emph{(étudier le quotient $v_{n+1}/v_n$)} ;
  \item $w_n=n+(-1)^{n}$ ;
  \item $t_n=\displaystyle\sum_{k=1}^{n}\frac1k$.
\end{questions}
\end{exercice}

\begin{exercice}[1]{Limites : opérations}
Déterminer les limites des suites suivantes :
\begin{tasks}(2)
  \task $\dfrac{3n^{2}+1}{n^{2}+2}$
  \task $\dfrac{2n+(-1)^{n}}{n}$
  \task $\sqrt{n+1}-\sqrt{n}$
  \task $\Bigl(1+\dfrac1n\Bigr)\Bigl(2-\dfrac{3}{n^{2}}\Bigr)$
  \task $n-n^{2}$
  \task $\dfrac{2^{n}}{n!}$\ \emph{(encadrer pour $n\geq2$)}
\end{tasks}
\end{exercice}

\begin{exercice}[1]{Théorème d'encadrement}
Déterminer les limites suivantes, à l'aide d'un encadrement :
\begin{questions}
  \item $\dfrac{n+\cos n}{n}$ ;
  \item $\dfrac{\sin n}{n^{2}}$ ;
  \item $\displaystyle\sum_{k=1}^{n}\frac{n}{n^{2}+k}$\quad
    \emph{(encadrer chaque terme de la somme)}.
\end{questions}
\end{exercice}

\begin{exercice}[1]{Arithmético-géométrique}
On définit $(u_n)$ par $u_0=0$ et, pour tout $n$,
$u_{n+1}=\dfrac{u_n}{2}+3$.
\begin{questions}
  \item Déterminer le réel $\omega$ tel que la suite constante égale à
    $\omega$ vérifie la récurrence.
  \item Montrer que $(u_n-\omega)$ est géométrique, puis exprimer $u_n$ en
    fonction de $n$.
  \item En déduire la limite de $(u_n)$.
\end{questions}
\end{exercice}

\begin{exercice}[1]{Avec la définition de la limite}
Revenir à la définition (avec $\varepsilon$ ou $A$) pour montrer :
\begin{questions}
  \item $\dfrac1n\longrightarrow0$ ;
  \item $\dfrac{n+1}{n}\longrightarrow1$ ;
  \item $n^{2}\longrightarrow+\infty$.
\end{questions}
\emph{Dans chaque cas, pour $\varepsilon>0$ (ou $A>0$) donné, on exhibera un
rang $N$ qui convient.}
\end{exercice}

\rubrique{Entraînement}

\begin{exercice}[2]{Opérations sur les bornes supérieures}
Soient $A$ et $B$ deux parties non vides et majorées de $\R$. On note
$A+B=\ens{a+b \mid a\in A,\ b\in B}$.
\begin{questions}
  \item Montrer que si $A\subset B$, alors $\sup A\leq\sup B$.
  \item Montrer que $\sup(A\cup B)=\max\bigl(\sup A,\sup B\bigr)$.
  \item Montrer que $\sup(A+B)=\sup A+\sup B$.
    \emph{Pour l'inégalité difficile, utiliser la caractérisation par
    $\varepsilon$.}
\end{questions}
\end{exercice}

\begin{exercice}[2]{Traduction séquentielle}
\begin{questions}
  \item Soit $X$ une partie non vide et majorée de $\R$, et $M=\sup X$.
    Montrer qu'il existe une suite $(x_n)$ d'éléments de $X$ telle que
    $x_n\to M$.
    \emph{Appliquer la caractérisation de la borne supérieure avec
    $\varepsilon=\frac{1}{n+1}$.}
  \item Pour $x\in\R$, on pose $d_n=\dfrac{\lfloor10^{n}x\rfloor}{10^{n}}$.
    Montrer que $(d_n)$ est une suite de décimaux qui converge vers $x$.
    Qu'en déduit-on pour $\mathbb{D}$ ?
\end{questions}
\end{exercice}

\begin{exercice}[2]{Premiers théorèmes, à la main}
\begin{questions}
  \item Montrer que si $u_n\to\ell>0$, alors $u_n>\dfrac{\ell}{2}$ à partir
    d'un certain rang.
  \item Montrer que si $(u_n)$ est bornée et $v_n\to0$, alors
    $u_n v_n\to0$.
  \item Montrer que si $u_n\to\ell$, alors $|u_n|\to|\ell|$. La réciproque
    est-elle vraie ?
\end{questions}
\end{exercice}

\begin{exercice}[2]{Le théorème de la limite monotone en action}
\begin{questions}
  \item On pose $u_n=\displaystyle\sum_{k=1}^{n}\frac{1}{k^{2}}$. Montrer que $(u_n)$
    est croissante et majorée \emph{(on rappelle l'inégalité
    $u_n\leq2-\frac1n$, démontrée au chapitre 1)}, et conclure.
  \item On pose $v_n=\displaystyle\sum_{k=0}^{n}\frac{1}{k!}$. Montrer que, pour
    $k\geq1$ : $\dfrac{1}{k!}\leq\dfrac{1}{2^{k-1}}$, en déduire que $(v_n)$
    est majorée par $3$, et conclure.
  \item Démontrer soigneusement : une suite croissante non majorée tend vers
    $+\infty$.
\end{questions}
\end{exercice}

\begin{exercice}[2]{Suites adjacentes}
\begin{questions}
  \item Vérifier que $a_n=1-\dfrac1n$ et $b_n=1+\dfrac1n$ ($n\geq1$) sont
    adjacentes. Quelle est leur limite commune ?
  \item On pose, pour $n\geq1$ :
    \[
      a_n=\sum_{k=0}^{n}\frac{1}{k!}
      \qquad\text{et}\qquad
      b_n=a_n+\frac{1}{n\cdot n!}.
    \]
    Montrer que $(a_n)$ et $(b_n)$ sont adjacentes. On note $\mathrm{e}$
    leur limite commune (c'est bien le nombre d'Euler) : donner
    l'encadrement de $\mathrm{e}$ fourni par $n=3$.
\end{questions}
\end{exercice}

\begin{exercice}[2]{Suites extraites}
\begin{questions}
  \item Montrer que $\bigl((-1)^{n}\bigr)$ diverge, à l'aide de deux suites
    extraites.
  \item La suite $\Bigl(\cos\dfrac{n\pi}{3}\Bigr)$ converge-t-elle ?
  \item Démontrer : si $u_{2n}\to\ell$ et $u_{2n+1}\to\ell$, alors
    $u_n\to\ell$.
  \item La suite définie par $u_n=(-1)^{n}+\dfrac1n$ converge-t-elle ?
\end{questions}
\end{exercice}

\begin{exercice}[2]{Récurrences linéaires d'ordre 2}
Déterminer le terme général de chacune des suites suivantes :
\begin{questions}
  \item $u_0=2$, $u_1=3$ et $u_{n+2}=3u_{n+1}-2u_n$ ;
  \item la suite de Fibonacci : $F_0=0$, $F_1=1$ et $F_{n+2}=F_{n+1}+F_n$
    \emph{(on fera apparaître le nombre d'or $\varphi=\frac{1+\sqrt5}{2}$)} ;
  \item $w_0=0$, $w_1=1$ et $w_{n+2}=w_{n+1}-w_n$. Que remarque-t-on ?
\end{questions}
\end{exercice}

\begin{exercice}[2]{Une suite récurrente complète}
On définit $(u_n)$ par $u_0=0$ et, pour tout $n\in\N$,
$u_{n+1}=\sqrt{2+u_n}$.
\begin{questions}
  \item Montrer que l'intervalle $\intFF{0}{2}$ est stable par
    $f:x\mapsto\sqrt{2+x}$, et en déduire que $0\leq u_n\leq2$ pour tout $n$.
  \item Étudier le signe de $f(x)-x$ sur $\intFF{0}{2}$, et en déduire que
    $(u_n)$ est croissante.
  \item Montrer que $(u_n)$ converge et déterminer sa limite.
  \item Illustrer la convergence par une construction en toile d'araignée.
\end{questions}
\end{exercice}

\begin{exercice}[2]{Une récurrence à deux visages}
On considère la récurrence $u_{n+1}=u_n^{2}$.
\begin{questions}
  \item Exprimer $u_n$ en fonction de $u_0$ et $n$.
  \item Étudier la convergence de $(u_n)$ selon la valeur de $u_0\geq0$
    (on distinguera $u_0<1$, $u_0=1$, $u_0>1$).
  \item Retrouver ces comportements par l'étude de $f:x\mapsto x^{2}$ :
    intervalles stables, signe de $f(x)-x$, points fixes.
\end{questions}
\end{exercice}

\begin{exercice}[2]{Suites complexes}
\begin{questions}
  \item Montrer que $z_n=\Bigl(\dfrac{1+\mathrm{i}}{2}\Bigr)^{n}$ converge
    vers $0$.
  \item On définit $(z_n)$ par $z_0=0$ et
    $z_{n+1}=\dfrac{\mathrm{i}}{2}\,z_n+1$. Déterminer le point fixe
    $\omega$, montrer que $|z_n-\omega|=|\omega|\cdot\Bigl(\dfrac12\Bigr)^{n}$,
    et en déduire la limite de $(z_n)$.
  \item Retrouver cette limite en séparant parties réelle et imaginaire pour
    les premiers termes : que valent $\operatorname{Re}\omega$ et
    $\operatorname{Im}\omega$ ?
\end{questions}
\end{exercice}

\rubrique{Approfondissement}

\begin{exercice}[3]{Le théorème de Cesàro}
Soit $(u_n)_{n\geq1}$ une suite convergeant vers $\ell$. On pose
$c_n=\dfrac{u_1+u_2+\dots+u_n}{n}$.
\begin{questions}
  \item Soit $\varepsilon>0$ et $N$ tel que $|u_k-\ell|\leq\varepsilon$ pour
    $k\geq N$. Montrer que, pour $n\geq N$ :
    \[
      |c_n-\ell|\leq\frac{1}{n}\sum_{k=1}^{N-1}|u_k-\ell|+\varepsilon.
    \]
  \item En déduire que $c_n\to\ell$ (théorème de Cesàro).
  \item Montrer, à l'aide de $u_n=(-1)^{n}$, que la réciproque est fausse.
  \item Application : montrer que si $u_{n+1}-u_n\to\ell$, alors
    $\dfrac{u_n}{n}\to\ell$.
\end{questions}
\end{exercice}

\begin{exercice}[3]{La série harmonique et la constante d'Euler}
Pour $n\geq1$, on pose $H_n=\displaystyle\sum_{k=1}^{n}\frac1k$.
\begin{questions}
  \item Montrer que, pour tout $k\geq1$ :
    $\dfrac{1}{k+1}\leq\ln(k+1)-\ln k\leq\dfrac1k$
    \emph{(utiliser l'encadrement $\frac{x}{1+x}\leq\ln(1+x)\leq x$ du
    chapitre 4)}.
  \item En sommant, montrer que $H_n\geq\ln(n+1)$, et en déduire que
    $H_n\to+\infty$.
  \item On pose $\gamma_n=H_n-\ln n$. Montrer que $(\gamma_n)$ est
    décroissante et minorée par $0$.
  \item Conclure : $(\gamma_n)$ converge. Sa limite $\gamma\approx0{,}577$
    est la \emph{constante d'Euler}.
\end{questions}
\end{exercice}

\begin{exercice}[3]{L'irrationalité de e}
On reprend les suites adjacentes de l'exercice 13 :
$a_n=\displaystyle\sum_{k=0}^{n}\frac{1}{k!}$ et $b_n=a_n+\dfrac{1}{n\cdot n!}$, de
limite commune $\mathrm{e}$, avec $a_n<\mathrm{e}<b_n$ pour tout $n\geq1$.
On suppose, par l'absurde, que $\mathrm{e}=\dfrac{p}{q}$ avec
$p,q\in\N^{*}$.
\begin{questions}
  \item Justifier que $q!\,a_q$ est un entier, puis que $q!\,\mathrm{e}$ est
    un entier.
  \item Montrer que $0<q!\,(\mathrm{e}-a_q)<\dfrac1q\leq1$.
  \item Conclure.
\end{questions}
\end{exercice}

\begin{exercice}[3]{La moyenne arithmético-géométrique}
Soient $0<b<a$. On définit deux suites par $a_0=a$, $b_0=b$ et :
\[
  a_{n+1}=\frac{a_n+b_n}{2},\qquad b_{n+1}=\sqrt{a_n\,b_n}.
\]
\begin{questions}
  \item Montrer que, pour tout $n$ : $0<b_n\leq a_n$
    \emph{(inégalité arithmético-géométrique, chapitre 2)}.
  \item Montrer que $(a_n)$ est décroissante et $(b_n)$ croissante.
  \item Montrer que $0\leq a_{n+1}-b_{n+1}\leq\dfrac{a_n-b_n}{2}$, puis que
    $a_n-b_n\to0$.
  \item Conclure que $(a_n)$ et $(b_n)$ convergent vers une même limite
    $M(a,b)$, la \emph{moyenne arithmético-géométrique} de $a$ et $b$, et
    donner l'encadrement obtenu pour $M(4,1)$ avec $n=2$.
\end{questions}
\end{exercice}

\begin{exercice}[3]{Bolzano-Weierstrass par dichotomie}
Soit $(u_n)$ une suite à valeurs dans un segment $\intFF{a}{b}$. On
construit deux suites $(a_n)$ et $(b_n)$ ainsi : $a_0=a$, $b_0=b$ ; si
l'intervalle $\Bigl[a_n,\frac{a_n+b_n}{2}\Bigr]$ contient $u_k$ pour une
infinité d'indices $k$, on le choisit comme $\intFF{a_{n+1}}{b_{n+1}}$ ;
sinon on choisit l'autre moitié.
\begin{questions}
  \item Justifier que, pour tout $n$, l'intervalle $\intFF{a_n}{b_n}$
    contient $u_k$ pour une infinité d'indices $k$.
  \item Montrer que $(a_n)$ et $(b_n)$ sont adjacentes ; on note $\ell$ leur
    limite commune.
  \item Construire une extraction $\varphi$ strictement croissante telle que
    $u_{\varphi(n)}\in\intFF{a_n}{b_n}$ pour tout $n$.
  \item Conclure : $u_{\varphi(n)}\to\ell$. Quel théorème vient-on de
    démontrer ?
\end{questions}
\end{exercice}

\begin{exercice}[3]{Fibonacci et le nombre d'or}
On reprend la suite de Fibonacci ($F_0=0$, $F_1=1$,
$F_{n+2}=F_{n+1}+F_n$) et son terme général obtenu à l'exercice 15 :
$F_n=\dfrac{\varphi^{n}-\psi^{n}}{\sqrt5}$, avec
$\varphi=\dfrac{1+\sqrt5}{2}$ et $\psi=\dfrac{1-\sqrt5}{2}$.
\begin{questions}
  \item Justifier que $|\psi|<1$, puis déterminer la limite de
    $\dfrac{F_{n+1}}{F_n}$.
  \item Démontrer l'identité de Cassini : pour tout $n\geq1$,
    $F_{n+1}F_{n-1}-F_n^{2}=(-1)^{n}$.
  \item On pose $r_n=\dfrac{F_{n+1}}{F_n}$ ($n\geq1$). Vérifier que
    $r_{n+1}=1+\dfrac{1}{r_n}$, et retrouver la limite de $(r_n)$ à l'aide
    de l'équation du point fixe.
  \item Expliquer pourquoi on peut écrire
    $\varphi=1+\cfrac{1}{1+\cfrac{1}{1+\cdots}}$ (fraction continuée).
\end{questions}
\end{exercice}

% !TEX root = ../../main.tex
%% Défi du chapitre 5 — La suite logistique, de la convergence au chaos
%% Résultats vérifiés par verifications/defis_01_05.py.

\rubrique{Défi}

\begin{defi}{La suite logistique, de la convergence au chaos}{2 h}
Pour $a\in\intOF{0}{4}$, on pose $f_a(x)=a\,x(1-x)$, et l'on considère la
suite définie par $u_0\in\intFF{0}{1}$ et $u_{n+1}=f_a(u_n)$. Ce modèle de
dynamique des populations, popularisé par le biologiste Robert May (1976),
présente selon la valeur de $a$ des comportements d'une étonnante
variété.

\partieExo{A}{Premières propriétés}
\begin{questions}
  \item Montrer que $f_a\bigl(\intFF{0}{1}\bigr)\subseteq\intFF{0}{1}$ ;
    en déduire que la suite est bien définie et à valeurs dans
    $\intFF{0}{1}$.
  \item Déterminer les points fixes de $f_a$ dans $\intFF{0}{1}$. On notera
    $\ell_a=1-\dfrac1a$ lorsque $a>1$.
  \item On suppose $a\in\intOF{0}{1}$. Montrer que $(u_n)$ est décroissante
    et converge vers $0$.
\end{questions}

\partieExo{B}{Convergence vers l'équilibre ($1<a\leqslant3$)}
\begin{questions}[resume]
  \item Calculer $f_a'(\ell_a)$ et vérifier que $\bigl|f_a'(\ell_a)\bigr|<1$
    pour $a\in\intOO{1}{3}$ (point fixe \emph{attractif}, chapitre 6).
  \item \emph{Le cas $a=2$.} Montrer que $1-2u_{n+1}=(1-2u_n)^{2}$ et en
    déduire une expression explicite de $u_n$. Pour $u_0\in\intOO{0}{1}$,
    étudier la convergence et sa vitesse.
  \item \emph{Le cas $1<a\leqslant2$.} Soit $u_0\in\intOO{0}{1}$. Montrer
    que $u_n\in\intOF{0}{\frac12}$ pour $n\geqslant1$, que $f_a$ est
    croissante sur cet intervalle, et en déduire que $(u_n)_{n\geqslant1}$
    est monotone et converge vers $\ell_a$.
\end{questions}

\medskip
\begin{center}
\begin{tikzpicture}[scale=5.2,line cap=round,line join=round]
  \draw[otRule,line width=0.3pt] (0,0) rectangle (1,1);
  \draw[otMuted,line width=0.5pt] (0,0) -- (1,1);
  \draw[otNavy,line width=1.1pt] plot[domain=0:1,samples=100] ({\x},{2.8*\x*(1-\x)});
  \draw[otGreenDark,line width=0.6pt]
    (0.0800,0.0000) -- (0.0800,0.2061) -- (0.2061,0.2061) -- (0.2061,0.4581) --
    (0.4581,0.4581) -- (0.4581,0.6951) -- (0.6951,0.6951) -- (0.6951,0.5934) --
    (0.5934,0.5934) -- (0.5934,0.6756) -- (0.6756,0.6756) -- (0.6756,0.6137) --
    (0.6137,0.6137) -- (0.6137,0.6638) -- (0.6638,0.6638) -- (0.6638,0.6249) --
    (0.6249,0.6249) -- (0.6249,0.6563) -- (0.6563,0.6563) -- (0.6563,0.6316) --
    (0.6316,0.6316) -- (0.6316,0.6515) -- (0.6515,0.6515) -- (0.6515,0.6357) --
    (0.6357,0.6357) -- (0.6357,0.6484) -- (0.6484,0.6484) -- (0.6484,0.6383);
  \fill[otGreenDark] (0.6429,0.6429) circle (0.25pt);
  \fill[otNavy] (0.08,0) circle (0.25pt);
  \node[font=\scriptsize,text=otNavy,below] at (0.08,0) {$u_0$};
  \node[font=\scriptsize,text=otMuted,below] at (1,0) {$1$};
  \node[font=\scriptsize,text=otMuted,left] at (0,1) {$1$};
  \node[font=\scriptsize,text=otNavy,anchor=west] at (0.88,0.42) {$y=f_a(x)$};
  \node[font=\scriptsize,text=otMuted,anchor=west] at (0.86,0.82) {$y=x$};
  \draw[otGreenDark,line width=0.4pt,dash pattern=on 1.5pt off 1.2pt] (0.6429,0.6429) -- (0.6429,0);
  \node[font=\scriptsize,text=otGreenDark,below] at (0.6429,0) {$\ell_a$};
\end{tikzpicture}\\[0.5mm]
\legende{La « toile d'araignée » pour $a=2{,}8$ : les itérés s'enroulent
autour du point fixe $\ell_a=\frac{9}{14}$, car $f_a'(\ell_a)=-0{,}8$ est
négatif.}
\end{center}

\partieExo{C}{Le chaos ($a=4$)}
\begin{questions}[resume]
  \item On écrit $u_0=\sin^{2}\theta$ avec $\theta\in\intFF{0}{\frac\pi2}$.
    Montrer que $u_n=\sin^{2}\bigl(2^{n}\theta\bigr)$ pour tout $n$.
  \item Montrer que pour $\theta=\dfrac{\pi}{5}$ la suite est périodique de
    période $2$, et calculer les deux valeurs prises ; montrer que pour
    $\theta=\dfrac{\pi}{7}$ elle est périodique de période $3$.
  \item Deux conditions initiales $\theta$ et $\theta+\varepsilon$ très
    proches donnent des angles $2^{n}\theta$ et $2^{n}(\theta+\varepsilon)$.
    À partir de quel rang les deux suites deviennent-elles sans rapport
    pour $\varepsilon=10^{-6}$ ? pour $\varepsilon=10^{-16}$ (précision
    d'un ordinateur) ? Commenter : c'est la \emph{sensibilité aux
    conditions initiales}, signature du chaos.
\end{questions}
\end{defi}

